IRS-assisted secure transmission, IRS-NOMA, and practical impairment modeling are all active topics, but the literature typically treats only a subset of those axes at once. In secure IRS-NOMA, Zhang \emph{et al.}~\cite{zhang2022tcom} optimize secrecy under outage constraints, while Wang \emph{et al.}~\cite{wang2022twc} strengthen an IRS-NOMA design with artificial jamming. STAR-RIS secure NOMA under imperfect eavesdropper CSI is studied by Jia \emph{et al.}~\cite{jia2023cl}, and heterogeneous-service coexistence for IRS-assisted uplink NOMA is examined by Na \emph{et al.}~\cite{na2024tvt}.

On the practical-modeling side, Shen \emph{et al.}~\cite{shen2021tcom} quantify the effect of transceiver hardware impairments on IRS beamforming, Yang \emph{et al.}~\cite{yang2023cl} analyze secure RIS systems under spatial correlation using statistical CSI, and Gao \emph{et al.}~\cite{gao2025tgcn} combine imperfect CSI, spatial correlation, and phase errors in an outage-constrained RIS design. These studies motivate the present work, but none of them simultaneously provide: 1) a secrecy-aware quantized IRS optimizer, 2) correlated Rician fading with CSI uncertainty and hardware loss in one simulation stack, and 3) a reproducible multi-metric package spanning secrecy, sum rate, energy efficiency, outage, fairness, and convergence.

Accordingly, this manuscript positions the proposed method as a robust secrecy-aware surrogate for practical IRS-NOMA operation. The comparison against AO is labeled as \emph{literature-inspired} rather than a claim of exact reproduction, while the overall novelty claim rests on the combined formulation, the robust projected solver, and the reproducible evidence stack generated directly from the codebase.
